% !TEX program = xelatex
% A题论文正文 — 问题四部分（按“同梦杯”论文格式整理）
\documentclass[UTF8,zihao=-4,a4paper]{ctexart}

\usepackage{amsmath,amssymb,amsfonts}
\usepackage{geometry}
\usepackage{booktabs}
\usepackage{array}
\usepackage{enumitem}
\usepackage{graphicx}
\usepackage{caption}
\usepackage{float}
\usepackage{hyperref}

\geometry{left=2.5cm,right=2.5cm,top=2.5cm,bottom=2.5cm}
\linespread{1.0}
\setlength{\parindent}{2em}
\setlength{\parskip}{0pt}
\setlist{nosep,leftmargin=2em}
\pagestyle{plain}

\ctexset{
  section={
    format=\centering\heiti\zihao{4},
    name={,、},
    number=\chinese{section},
    aftername=\quad,
    beforeskip=1.0ex,
    afterskip=1.0ex
  },
  subsection={
    format=\raggedright\heiti\zihao{-4},
    name={,},
    number=\chinese{section}.\arabic{subsection},
    aftername=\quad,
    beforeskip=0.8ex,
    afterskip=0.6ex
  },
  subsubsection={
    format=\raggedright\heiti\zihao{-4},
    name={,},
    number=\chinese{section}.\arabic{subsection}.\arabic{subsubsection},
    aftername=\quad,
    beforeskip=0.6ex,
    afterskip=0.4ex
  }
}

\captionsetup{font=small,labelfont=bf}
\renewcommand{\arraystretch}{1.15}

\begin{document}

\begin{center}
{\heiti\zihao{3} 面向管理决策的校园微交通综合治理建议}
\end{center}

\section{问题重述}

问题四要求基于前三问的建模结果，用通俗、清晰且具有管理可执行性的语言，为浙师大管理层撰写一页纸建议书。该建议书需要说明校园能够安全容纳多少非机动车，并说服管理层采用前文设计的路网优化、停车治理和微公交接驳方案。因此，问题四并不是重新建立一个独立的技术模型，而是将问题一至问题三的模型结论转化为管理层能够理解和执行的治理方案。

从前三问结果看，校园微交通治理的核心矛盾包括三个方面：第一，高峰期主干道人车混行导致冲突风险较高；第二，电动车实际保有量已经超过校园安全环境承载力，且16号楼、25号楼、商业街等核心区域停车空间严重不足；第三，微公交能够替代部分高峰私人电动车直达需求，但必须与核心区限停、外围分流和路网安全优化配套实施，才能形成稳定治理效果。

\section{问题分析}

\subsection{问题四的分析}

问题四的写作对象是学校管理层，因此重点应从“模型推导”转向“结果解释”和“政策建议”。前三问已经分别回答了路网怎么优化、电动车能容纳多少、微公交如何规划和调度。本问需要把这些结果串联成一个完整的治理逻辑，即：先明确校园非机动车安全承载力，再说明现有交通风险和停车压力，最后提出分阶段、可落地的综合治理方案。

由于管理建议需要排序，本文在前三问模型结果基础上增加一个轻量级治理优先级评价框架。该框架不作为新的复杂数学模型，而用于将安全改善、容量缓解和实施可行性三个维度统一起来，帮助管理层判断哪些措施应优先推进。这样既能保持建议书的通俗性，又能体现建议来源于定量模型而非主观判断。

\section{模型假设}

为便于将前三问结果转化为管理建议，作如下假设：
\begin{enumerate}
    \item 学校管理目标以安全优先为基本原则，在安全风险明显下降的前提下，可以接受小幅通行时间增加；
    \item 校园电动车治理既要考虑高峰同时运行需求，也要考虑实际保有量形成的长期存量压力；
    \item 微公交主要用于替代高峰私人电动车直达需求，不能单独吸收全部存量电动车；
    \item 管理建议应优先选择安全收益明显、容量缓解效果较好且短期内具备实施条件的措施；
    \item 第四问建议书面向管理层，表达应尽量避免复杂公式，突出关键数字、治理原因和实施路径。
\end{enumerate}

\section{符号说明}

\begin{table}[H]
\centering
\caption{问题四主要符号说明}
\label{tab:symbols_q4}
\begin{tabular}{ccc}
\toprule
\textbf{符号} & \textbf{说明} & \textbf{单位} \\
\midrule
$N_{max}$ & 校园电动车安全环境承载力 & 辆 \\
$N_{stock}$ & 校园电动车实际保有量估计值 & 辆 \\
$\theta_{peak}$ & 高峰运行饱和度 & -- \\
$\theta_{stock}$ & 存量饱和度 & -- \\
$TTT$ & 全路网总行程时间 & 人$\cdot$s \\
$S$ & 人车冲突风险指标 & -- \\
$D_s^t$ & 时段 $t$ 微公交接驳需求 & 人次 \\
$Priority_k$ & 治理措施 $k$ 的综合优先级得分 & -- \\
\bottomrule
\end{tabular}
\end{table}

\section{模型的建立与求解}

\subsection{问题四结果转化与建议书形成}

\subsubsection{前三问核心结果汇总}

问题一建立了校园多方式交通流分配与分流优化模型。结果表明，综合优化方案C能够使人车冲突风险下降22.0\%，但总行程时间增加2.1\%。因此，方案C不应表述为单纯提升通行效率的方案，而应定位为“安全优先型”路网优化方案。该结果说明，在校园交通场景下，降低人车混行冲突比追求最短通行时间更重要。

问题二建立了电动车环境承载力评估模型。扣除消防通道、建筑出入口、人行通道和必要疏散空间后，校园电动车安全环境承载力为
\begin{equation}
N_{max}=7627\text{ 辆}.
\end{equation}
按高峰运行OD需求校核，最差时段电动车运行需求为5323辆，对应高峰运行饱和度为
\begin{equation}
\theta_{peak}=\frac{5323}{7627}=0.698.
\end{equation}
这说明高峰期并非所有车辆同时进入道路和停车系统。但若按公开资料中15000辆以上的实际保有量进行存量校核，则
\begin{equation}
\theta_{stock}=\frac{15000}{7627}\approx1.97,
\end{equation}
安全容量缺口约为
\begin{equation}
15000-7627=7373\text{ 辆}.
\end{equation}
因此，校园电动车治理不能简单理解为“高峰不超载即可”，而应认识到实际保有规模已经明显超过安全承载力。

问题三建立了微公交站点选址、线路规划和潮汐调度模型。推荐接驳线路为
\begin{equation}
\text{桃园站}\rightarrow\text{16号楼站}\rightarrow\text{25号楼站}\rightarrow\text{杏园站}\rightarrow\text{桂苑站}.
\end{equation}
基准10\%替代率下，早高峰、午高峰和晚高峰接驳需求分别为415、439和532人次，对应投放车辆数分别为3辆、3辆和4辆。接驳实施后，路网总行程时间改善约2.6\%--2.9\%，人车冲突风险改善约6.6\%--7.0\%。

\subsubsection{治理措施优先级评价框架}

为将前三问结果转化为管理建议，设治理措施 $k$ 的综合优先级得分为
\begin{equation}
Priority_k=0.40Safety_k+0.35Capacity_k+0.25Feasibility_k,
\end{equation}
其中 $Safety_k$ 表示安全改善程度，$Capacity_k$ 表示容量缓解程度，$Feasibility_k$ 表示实施可行性。权重设置体现校园治理中安全优先、兼顾容量和可操作性的原则。由于第4问面向管理建议，该评分仅用于措施排序，不再展开复杂参数标定。

\begin{table}[H]
\centering
\caption{校园微交通治理措施优先级排序}
\label{tab:priority_q4}
\begin{tabular}{ccccc}
\toprule
\textbf{治理措施} & \textbf{安全改善} & \textbf{容量缓解} & \textbf{可行性} & \textbf{优先级} \\
\midrule
清理消防通道、划定停车边界 & 高 & 高 & 高 & 第一优先 \\
16号楼、25号楼、商业街核心区限停 & 高 & 高 & 中 & 第一优先 \\
开通教学潮汐微公交线路 & 中 & 中 & 高 & 第二优先 \\
主干道人车隔离与分时限行 & 高 & 中 & 中 & 第二优先 \\
注册准入与存量车辆分期管控 & 高 & 高 & 中 & 长期推进 \\
外围集中停放与步行接驳引导 & 中 & 高 & 中 & 长期推进 \\
\bottomrule
\end{tabular}
\end{table}

由表\ref{tab:priority_q4}可知，短期内最应优先推进的是消防通道清理、停车边界划定和核心区限停。这些措施实施成本较低，但能直接改善乱停放和安全隐患。中期应开通教学潮汐微公交线路，并同步推进主干道人车隔离与分时限行。长期则应结合注册准入和存量车辆分期管控，将电动车规模逐步压降到安全承载力附近。

\subsubsection{面向管理层的一页纸建议书}

\begin{center}
{\heiti 关于浙师大金华校区校园“微交通”综合治理的建议书}
\end{center}

根据对校园道路、人车混行、电动车停放和微公交接驳系统的综合建模分析，我们建议学校采用“安全优先、总量控制、核心区限停、外围分流、微公交替代”的综合治理方案。

首先，校园电动车规模已经超过安全承载边界。模型测算表明，在严格扣除消防通道、建筑出入口、人行通道和必要疏散空间后，金华校区电动车安全环境承载力约为7627辆。虽然高峰时段模型估计的同时运行需求约为5323辆，尚未超过安全承载力，但公开资料显示校园电动车保有量已超过15000辆，约为安全承载力的1.97倍，缺口约7373辆。这说明学校当前面临的主要问题不是某一个时段车辆同时出行过多，而是长期车辆存量已经明显超过校园安全停放和管理能力。

其次，校园交通治理应坚持安全优先。问题一结果显示，综合路网优化方案虽然会使总行程时间增加2.1\%，但可使人车冲突风险下降22.0\%。因此，建议学校在16号楼、25号楼、食堂和商业街周边重点推行人车隔离、分时限行和局部单向组织。对于高峰期人流、电动车流交织严重的主干道，应优先设置清晰的通行边界，减少行人、电动车和机动车之间的交叉冲突。

第三，应立即治理核心区停车压力。16号楼、25号楼和商业街周边停车饱和度明显偏高，是校园停车矛盾最突出的区域。建议学校将这些区域设为核心限停区，严格禁止占用消防通道、建筑出入口和主通行人行道；同时将东区、西区新教学楼周边容量相对充足的区域作为外围缓冲停车区，引导师生“外围停车、步行或接驳进入核心区”。

第四，建议开通教学潮汐微公交线路。推荐线路为“桃园站—16号楼站—25号楼站—杏园站—桂苑站”。该线路连接主要生活区和核心教学区，能够服务早高峰宿舍到教学楼、午高峰教学楼到食堂及生活区、晚高峰教学楼返回宿舍的潮汐需求。基准方案下，早高峰、午高峰、晚高峰分别投放3辆、3辆、4辆接驳车即可满足连续发车要求。模型估计，该线路可使总行程时间下降约2.6\%--2.9\%，人车冲突风险下降约6.6\%--7.0\%。需要说明的是，微公交不能一次性替代全部存量电动车，但可以有效减少高峰期私人电动车直达核心区的需求。

最后，建议学校分阶段推进治理。第一阶段，立即清理消防通道、划定停车边界，并对16号楼、25号楼、商业街实行核心区限停；第二阶段，开通教学潮汐微公交线路，并配套设置站点标识、候车区和运行时刻表；第三阶段，建立电动车注册准入和分区管理制度，以7627辆安全承载力作为长期总量控制参考，逐步压降超出安全承载力的存量车辆。通过上述措施，学校可以在不大规模改造校园道路的前提下，有效降低人车冲突风险，缓解核心区停车压力，形成更加安全、有序、绿色的校园微交通系统。

\subsubsection{问题四结论}

综合前三问结果，校园微交通治理应从单点治理转向系统治理。问题一说明路网优化的重点是降低冲突风险，问题二说明校园实际电动车存量已明显超过安全承载力，问题三说明微公交能够承担部分高峰替代和核心区分流功能。因此，第4问建议不再建立复杂新模型，而是将前三问结果转化为管理建议：以7627辆作为安全承载力参考上限，以核心区限停和消防通道清理作为近期重点，以教学潮汐微公交作为中期替代工具，以注册准入和存量管控作为长期治理手段。该方案具有明确的数据依据和较强的可执行性。

\end{document}
